
\chapter{2021年北京卷}

\section{单选题}

\begin{problem}
已知集合 \(A=\{x\mid -1<x<1\}\)，\(B=\{x\mid 0\le x\le 2\}\)，则 \(A\cup B=\)（\quad）

\choices
    {\(( -1,2 )\)}
    {\(( -1,2 ]\)}
    {\([0,1)\)}
    {\([0,1]\)}

\end{problem}

\begin{answer}
B
\end{answer}

\begin{solution}
由题意可得：\(A\cup B=\{x\mid -1<x\le 2\}\)，即 \(A\cup B=( -1,2 ]\). 故选：B.
\end{solution}

\begin{problem}
在复平面内，复数 \(z\) 满足 \((1-\i)z=2\)，则 \(z=\)（\quad）

\choices
    {\(2+\i\)}
    {\(2-\i\)}
    {\(1-\i\)}
    {\(1+\i\)}

\end{problem}

\begin{answer}
D
\end{answer}

\begin{solution}
由题意可得 \(z=\frac{2}{1-\i}=\frac{2(1+\i)}{(1-\i)(1+\i)}=1+\i\). 故选：D.
\end{solution}

\begin{problem}
已知函数 \(f(x)\) 是定义在 \([0,1]\) 上的函数，那么“函数 \(f(x)\) 在 \([0,1]\) 上单调递增”是“函数 \(f(x)\) 在 \([0,1]\) 上的最大值为 \(f(1)\)”的（\quad）

\choices
    {充分而不必要条件}
    {必要而不充分条件}
    {充分必要条件}
    {既不充分也不必要条件}

\end{problem}

\begin{answer}
A
\end{answer}

\begin{solution}
若函数 \(f(x)\) 在 \([0,1]\) 上单调递增，则 \(f(x)\) 在 \([0,1]\) 上的最大值为 \(f(1)\).

若函数 \(f(x)\) 在 \([0,1]\) 上的最大值为 \(f(1)\)，比如 \(f(x)=\left( x-\frac{1}{3} \right)^2\)，但 \(f(x)=\left( x-\frac{1}{3} \right)^2\) 在 \(\left[0,\frac{1}{3}\right]\) 为减函数，在 \(\left[\frac{1}{3},1\right]\) 为增函数，故 \(f(x)\) 在 \([0,1]\) 上的最大值为 \(f(1)\) 推不出 \(f(x)\) 在 \([0,1]\) 上单调递增.

故“函数 \(f(x)\) 在 \([0,1]\) 上单调递增”是“函数 \(f(x)\) 在 \([0,1]\) 上的最大值为 \(f(1)\)”的充分而不必要条件. 故选：A.
\end{solution}

\begin{problem}
某四面体的三视图如图所示，该四面体的表面积为（\quad）

\bitmapfigure[max width=6.0cm,max totalheight=4.8cm]{img/2021/beijing/q4_fig1.png}

\choices
    {\(\frac{3+\sqrt3}{2}\)}
    {4}
    {\(3+\sqrt3\)}
    {2}

\end{problem}

\begin{answer}
A
\end{answer}

\begin{solution}
根据三视图可得如图所示的几何体，为正三棱锥 \(O-ABC\)，其侧面为等腰直角三角形，底面为等边三角形.

\bitmapfigure[max width=4.6cm,max totalheight=4.8cm]{img/2021/beijing/q4_solution_fig1.png}

由三视图可得该正三棱锥的侧棱长为 \(1\)，故其表面积为 \(3\times\frac{1}{2}\times1\times1+\frac{\sqrt3}{4}\times\left( \sqrt2 \right)^2=\frac{3+\sqrt3}{2}\).

故选：A.
\end{solution}

\begin{problem}
双曲线 \(C:\frac{x^2}{a^2}-\frac{y^2}{b^2}=1\) 过点 \(( \sqrt2,\sqrt3 )\)，且离心率为 \(2\)，则该双曲线的标准方程为（\quad）

\choices
    {\(x^2-\frac{y^2}{3}=1\)}
    {\(\frac{x^2}{3}-y^2=1\)}
    {\(x^2-\frac{\sqrt3 y^2}{3}=1\)}
    {\(\frac{\sqrt3 x^2}{3}-y^2=1\)}

\end{problem}

\begin{answer}
A
\end{answer}

\begin{solution}
由 \(e=\frac{c}{a}=2\)，可得 \(c=2a\)，\(b=\sqrt{c^2-a^2}=\sqrt3 a\).
所以双曲线的方程为 \(\frac{x^2}{a^2}-\frac{y^2}{3a^2}=1\).
将点 \(( \sqrt2,\sqrt3 )\) 的坐标代入双曲线的方程可得 \(\frac{2}{a^2}-\frac{3}{3a^2}=1\)，解得 \(a=1\)，故 \(b=\sqrt3\).
因此，双曲线的方程为 \(x^2-\frac{y^2}{3}=1\).
故选：A.
\end{solution}

\begin{problem}
\(\{a_n\}\) 和 \(\{b_n\}\) 是两个等差数列，其中 \(\frac{a_k}{b_k}\)（\(1\le k\le5\)）为常值，\(a_1=288\)，\(a_5=96\)，\(b_1=192\)，则 \(b_3=\)（\quad）

\choices
    {64}
    {128}
    {256}
    {512}

\end{problem}

\begin{answer}
B
\end{answer}

\begin{solution}
由已知条件可得 \(\frac{a_1}{b_1}=\frac{a_5}{b_5}\)，则 \(b_5=\frac{a_5b_1}{a_1}=\frac{96\times192}{288}=64\).
因此，\(b_3=\frac{b_1+b_5}{2}=\frac{192+64}{2}=128\).
故选：B.
\end{solution}

\begin{problem}
函数 \(f(x)=\cos x-\cos 2x\)，试判断函数的奇偶性及最大值（\quad）

\choices
    {奇函数，最大值为 \(2\)}
    {偶函数，最大值为 \(2\)}
    {奇函数，最大值为 \(\frac{9}{8}\)}
    {偶函数，最大值为 \(\frac{9}{8}\)}

\end{problem}

\begin{answer}
D
\end{answer}

\begin{solution}
由题意，\(f(-x)=\cos(-x)-\cos(-2x)=\cos x-\cos 2x=f(x)\)，所以该函数为偶函数.

又 \(f(x)=\cos x-\cos 2x=-2\cos^2 x+\cos x+1=-2\left( \cos x-\frac{1}{4} \right)^2+\frac{9}{8}\)，所以当 \(\cos x=\frac{1}{4}\) 时，\(f(x)\) 取最大值 \(\frac{9}{8}\).

故选：D.
\end{solution}

\begin{problem}
定义：24 小时内降水在平地上积水厚度（mm）来判断降雨程度，其中小雨（\(<10\text{ mm}\)），中雨（\(10\text{ mm}-25\text{ mm}\)），大雨（\(25\text{ mm}-50\text{ mm}\)），暴雨（\(50\text{ mm}-100\text{ mm}\)），小明用一个圆锥形容器接了 24 小时的雨水，如图，则这天降雨属于哪个等级（\quad）

\bitmapfigure[max width=6.0cm,max totalheight=4.8cm]{img/2021/beijing/q8_fig1.png}

\choices
    {小雨}
    {中雨}
    {大雨}
    {暴雨}

\end{problem}

\begin{answer}
B
\end{answer}

\begin{solution}
由题意，一个半径为 \(\frac{200}{2}=100\text{ mm}\) 的圆面内的降雨充满一个底面半径为 \(\frac{200}{2}\times\frac{150}{300}=50\text{ mm}\)，高为 \(150\text{ mm}\) 的圆锥.
所以积水厚度 \(d=\frac{\frac{1}{3}\pi\times50^2\times150}{\pi\times100^2}=12.5\text{ mm}\)，属于中雨.
故选：B.
\end{solution}

\begin{problem}
已知圆 \(C:x^2+y^2=4\)，直线 \(l:y=kx+m\)，当 \(k\) 变化时，\(l\) 截得圆 \(C\) 弦长的最小值为 \(2\)，则 \(m=\)（\quad）

\choices
    {\(\pm2\)}
    {\(\pm\sqrt2\)}
    {\(\pm\sqrt3\)}
    {\(\pm\sqrt5\)}

\end{problem}

\begin{answer}
C
\end{answer}

\begin{solution}
由题可得圆心为 \((0,0)\)，半径为 \(2\)，则圆心到直线的距离 \(d=\frac{|m|}{\sqrt{k^2+1}}\).
故弦长为 \(2\sqrt{4-\frac{m^2}{k^2+1}}\).
当 \(k=0\) 时，弦长取得最小值 \(2\sqrt{4-m^2}=2\)，解得 \(m=\pm\sqrt3\).
故选：C.
\end{solution}

\begin{problem}
数列 \(\{a_n\}\) 是递增的整数数列，且 \(a_1\ge3\)，\(a_1+a_2+\cdots+a_n=100\)，则 \(n\) 的最大值为（\quad）

\choices
    {9}
    {10}
    {11}
    {12}

\end{problem}

\begin{answer}
C
\end{answer}

\begin{solution}
若要使 \(n\) 尽可能大，则 \(a_1\)，递增幅度要尽可能小.
不妨设数列 \(\{a_n\}\) 是首项为 \(3\)，公差为 \(1\) 的等差数列，其前 \(n\) 项和为 \(S_n\)，则 \(a_n=n+2\)，\(S_{11}=\frac{3+13}{2}\times11=88<100\)，\(S_{12}=\frac{3+14}{2}\times12=102>100\).
所以 \(n\) 的最大值为 \(11\).
故选：C.
\end{solution}
\section{填空题}

\begin{problem}
在 \(\left( x^3-\frac{1}{x} \right)^4\) 展开式中常数项为\(\fillinblank{}\).

\end{problem}

\begin{answer}
\(-4\)
\end{answer}

\begin{solution}
试题分析：
\(\left( x^3-\frac{1}{x} \right)^4\)
的展开式的通项
\(T_{r+1}=C_4^r\left( x^3 \right)^{4-r}\left( -\frac{1}{x} \right)^r=\left( -1 \right)^r C_4^r x^{12-4r}\)，
令 \(r=3\) 得常数项为
\(T_4=\left( -1 \right)^3C_4^3=-4\).
\end{solution}

\begin{problem}
已知抛物线 \(C:y^2=4x\)，焦点为 \(F\)，点 \(M\) 为抛物线 \(C\) 上的点，且 \(\left| MF \right|=6\)，则 \(M\) 的横坐标是 \(\fillinblank{}\)；作 \(MN\perp x\) 轴于 \(N\)，则 \(S_{\triangle FMN}=\fillinblank{}\).

\end{problem}

\begin{answer}
\(5,\ 4\sqrt5\)
\end{answer}

\begin{solution}
因为抛物线的方程为 \(y^2=4x\)，故 \(p=2\) 且 \(F(1,0)\).
因为 \(\left| MF \right|=6\)，\(x_M+\frac{p}{2}=6\)，解得 \(x_M=5\)，故 \(y_M=\pm2\sqrt5\).
所以 \(S_{\triangle FMN}=\frac{1}{2}\times( 5-1 )\times2\sqrt5=4\sqrt5\).
故答案为：\(5\)，\(4\sqrt5\).
\end{solution}

\begin{problem}
已知 \(\bs{a}=(2,1)\)，\(\bs{b}=(2,-1)\)，\(\bs{c}=(0,1)\)，则 \(( \bs{a}+\bs{b} )\cdot\bs{c}\) 的值是 \(\fillinblank{}\)；\(\bs{a}\cdot\bs{b}\) 的值是 \(\fillinblank{}\).

\end{problem}

\begin{answer}
\(0,\ 3\)
\end{answer}

\begin{solution}
\(\bs{a}=(2,1)\)，\(\bs{b}=(2,-1)\)，\(\bs{c}=(0,1)\)，\((\bs{a}+\bs{b})=(4,0)\)，故 \((\bs{a}+\bs{b})\cdot\bs{c}=4\times0+0\times1=0\).
又 \(\bs{a}\cdot\bs{b}=2\times2+1\times( -1 )=3\).

故答案为：\(0\)；\(3\).
\end{solution}

\begin{problem}
若点 \(P(\cos\theta,\sin\theta)\) 与点 \(Q\left( \cos\left( \theta+\frac{\pi}{6} \right),\sin\left( \theta+\frac{\pi}{6} \right) \right)\) 关于 \(y\) 轴对称，写出一个符合题意的 \(\theta=\fillinblank{}\).

\end{problem}

\begin{answer}
\(\frac{5\pi}{12}\)
\end{answer}

\begin{solution}
根据 \(P\)，\(Q\) 在单位圆上，可得 \(\theta\)，\(\theta+\frac{\pi}{6}\) 关于 \(y\) 轴对称，得出 \(\theta+\frac{\pi}{6}+\theta=\pi+2k\pi\)，\(k\in\Z\).
解得 \(\theta=k\pi+\frac{5\pi}{12}\)，\(k\in\Z\).
当 \(k=0\) 时，可取 \(\theta\) 的一个值为 \(\frac{5\pi}{12}\).
故答案为：\(\frac{5\pi}{12}\)（满足 \(\theta=k\pi+\frac{5\pi}{12}\)，\(k\in\Z\) 即可）.
\end{solution}

\begin{problem}
已知函数 \(f(x)=\left| \lg x \right|-kx-2\)，给出下列四个结论：

\begin{circlelist}
\item 若 \(k=0\)，则 \(f(x)\) 有两个零点；

\item \(\exists k<0\)，使得 \(f(x)\) 有一个零点；

\item \(\exists k<0\)，使得 \(f(x)\) 有三个零点；

\item \(\exists k>0\)，使得 \(f(x)\) 有三个零点.
\end{circlelist}

以上正确结论的序号是\(\fillinblank{}\).

\end{problem}

\begin{answer}
\(\circled{1}\circled{2}\circled{4}\)
\end{answer}

\begin{solution}
由 \(f(x)=0\) 可得出 \(\left| \lg x \right|=kx+2\).
考查直线 \(y=kx+2\) 与曲线 \(g(x)=\left| \lg x \right|\) 的左，右支分别相切的情形，利用方程思想以及数形结合可判断各选项的正误.

对于 \(\circled{1}\)，当 \(k=0\) 时，由 \(f(x)=\left| \lg x \right|-2=0\)，
可得 \(x=\frac{1}{100}\) 或 \(x=100\)，\(\circled{1}\) 正确.

对于 \(\circled{2}\)，考查直线 \(y=kx+2\) 与曲线 \(y=-\lg x\)（\(0<x<1\)）相切于点 \(P(t,-\lg t)\). 对函数 \(y=-\lg x\) 求导得 \(y'=-\frac{1}{x\ln10}\)，由题意可得 \(kt+2=-\lg t\)，\(k=-\frac{1}{t\ln10}\).
解得 \(t=\frac{\e}{100}\)，\(k=-\frac{100}{\e}\lg\e\).
所以，存在 \(k=-\frac{100}{\e}\lg\e<0\)，
使得 \(f(x)\) 只有一个零点，\(\circled{2}\) 正确.

对于 \(\circled{3}\)，当直线 \(y=kx+2\) 过点 \((1,0)\) 时，\(k+2=0\)，解得 \(k=-2\). 所以，当 \(-\frac{100}{\e}\lg\e<k<-2\) 时，直线 \(y=kx+2\) 与曲线 \(y=-\lg x\)（\(0<x<1\)）有两个交点. 若函数 \(f(x)\) 有三个零点，则直线 \(y=kx+2\) 与曲线 \(y=-\lg x\)（\(0<x<1\)）有两个交点，直线 \(y=kx+2\) 与曲线 \(y=\lg x\)（\(x>1\)）有一个交点，所以 \(-\frac{100}{\e}\lg\e<k<-2\) 且 \(k+2>0\)，此不等式无解，因此，不存在 \(k<0\)，使得函数 \(f(x)\) 有三个零点，\(\circled{3}\) 错误.

对于 \(\circled{4}\)，考查直线 \(y=kx+2\) 与曲线 \(y=\lg x\)（\(x>1\)）相切于点 \(P(t,\lg t)\). 对函数 \(y=\lg x\) 求导得 \(y'=\frac{1}{x\ln10}\). 由题意可得 \(kt+2=\lg t\)，\(k=\frac{1}{t\ln10}\).
解得 \(t=100\e\)，\(k=\frac{\lg\e}{100\e}\).
所以，当 \(0<k<\frac{\lg\e}{100\e}\) 时，函数 \(f(x)\) 有三个零点，\(\circled{4}\) 正确.

故答案为：\(\circled{1}\circled{2}\circled{4}\).
\end{solution}
\section{解答题}

\begin{problem}
已知在三角形 \(\triangle ABC\) 中，\(c=2b\cos B\)，\(C=\frac{2\pi}{3}\).

\begin{enumerate}
    \item 求 \(B\) 的大小；
    \item 在下列三个条件中选择一个作为已知，使 \(\triangle ABC\) 存在且唯一确定，并求出 \(BC\) 边上的中线的长度.

    条件
\begin{circlelist}
\item ：\(c=\sqrt2 b\)；条件
\item ：周长为 \(4+2\sqrt3\)；条件
\item ：面积为 \(S_{\triangle ABC}=\frac{3\sqrt3}{4}\).
\end{circlelist}
\end{enumerate}

\end{problem}

\begin{solution}
（1）由 \(c=2b\cos B\)，则由正弦定理可得 \(\sin C=2\sin B\cos B\).
所以 \(\sin2B=\sin\frac{2\pi}{3}=\frac{\sqrt3}{2}\).
又因为 \(C=\frac{2\pi}{3}\)，所以 \(B\in\left( 0,\frac{\pi}{3} \right)\)，\(2B\in\left( 0,\frac{2\pi}{3} \right)\)，故 \(2B=\frac{\pi}{3}\)，解得 \(B=\frac{\pi}{6}\).

（2）若选择条件 \circled{1}，由正弦定理结合（1）可得
\(\frac{c}{b}=\frac{\sin C}{\sin B}=\frac{\frac{\sqrt3}{2}}{\frac{1}{2}}=\sqrt3\)，与 \(c=\sqrt2 b\) 矛盾，故这样的 \(\triangle ABC\) 不存在.

若选择条件 \circled{2}，由（1）可得 \(A=\frac{\pi}{6}\).
设 \(\triangle ABC\) 的外接圆半径为 \(R\)，则由正弦定理可得 \(a=b=2R\sin\frac{\pi}{6}=R\)，\(c=2R\sin\frac{2\pi}{3}=\sqrt3 R\).
于是周长 \(a+b+c=2R+\sqrt3 R=4+2\sqrt3\)，解得 \(R=2\)，故 \(a=2\)，\(c=2\sqrt3\).
由余弦定理可得 \(BC\) 边上的中线的长度为 \(\sqrt{\left( 2\sqrt3 \right)^2+1^2-2\times2\sqrt3\times1\times\cos\frac{\pi}{6}}=\sqrt7\).

若选择条件 \circled{3}，由（1）可得 \(A=\frac{\pi}{6}\)，即 \(a=b\). 又 \(S_{\triangle ABC}=\frac{1}{2}ab\sin C=\frac{1}{2}a^2\times\frac{\sqrt3}{2}=\frac{3\sqrt3}{4}\)，解得 \(a=\sqrt3\).
由余弦定理可得 \(BC\) 边上的中线的长度为 \(\sqrt{b^2+\left( \frac{a}{2} \right)^2-2\times b\times\frac{a}{2}\times\cos\frac{2\pi}{3}}=\sqrt{3+\frac{3}{4}+\sqrt3\times\frac{\sqrt3}{2}}=\frac{\sqrt{21}}{2}\).
\end{solution}

\begin{problem}
已知正方体 \(ABCD-A_1B_1C_1D_1\)，点 \(E\) 为 \(A_1D_1\) 中点，直线 \(B_1C_1\) 交平面 \(CDE\) 于点 \(F\).

\begin{enumerate}
    \item 证明：点 \(F\) 为 \(B_1C_1\) 的中点；
    \item 若点 \(M\) 为棱 \(A_1B_1\) 上一点，且二面角 \(M-CF-E\) 的余弦值为 \(\frac{\sqrt5}{3}\)，求 \(\frac{A_1M}{A_1B_1}\) 的值.
\end{enumerate}

\bitmapfigure[max width=6.0cm,max totalheight=4.8cm]{img/2021/beijing/q17_fig1.png}

\end{problem}

\begin{solution}
（1）如图所示，取 \(B_1C_1\) 的中点 \(F'\)，连结 \(DE\)，\(EF'\)，\(F'C\).

\bitmapfigure[max width=4.8cm,max totalheight=4.8cm]{img/2021/beijing/q17_solution_fig1.png}

由于 \(ABCD-A_1B_1C_1D_1\) 为正方体，\(E\)，\(F'\) 为中点，故 \(EF'//CD\).

从而 \(E\)，\(F'\)，\(C\)，\(D\) 四点共面，即平面 \(CDE\) 即平面 \(CDEF'\).

据此可得：直线 \(B_1C_1\) 交平面 \(CDE\) 于点 \(F'\).

当直线与平面相交时只有唯一的交点，故点 \(F\) 与点 \(F'\) 重合，即点 \(F\) 为 \(B_1C_1\) 中点.

（2）以点 \(D\) 为坐标原点，\(DA\)，\(DC\)，\(DD_1\) 方向分别为 \(x\) 轴，\(y\) 轴，\(z\) 轴正方向，建立空间直角坐标系 \(D-xyz\).

\bitmapfigure[max width=4.8cm,max totalheight=4.8cm]{img/2021/beijing/q17_solution_fig2.png}

不妨设正方体的棱长为 \(2\)，设
 \(\frac{A_1M}{A_1B_1}=\lambda\)（\(0\le\lambda\le1\)）.
则
 \(M(2,2\lambda,2)\)，\(C(0,2,0)\)，\(F(1,2,2)\)，\(E(1,0,2)\).
从而
 \(\overrightarrow{MC}=(-2,2-2\lambda,-2)\)，\(\overrightarrow{CF}=(1,0,2)\)，\(\overrightarrow{FE}=(0,-2,0)\).

设平面 \(MCF\) 的法向量为 \(\bs{m}=(x_1,y_1,z_1)\)，则
\examdisplaycases{}{
\bs{m}\cdot\overrightarrow{MC}=-2x_1+(2-2\lambda)y_1-2z_1=0,\\
\bs{m}\cdot\overrightarrow{CF}=x_1+2z_1=0.
}
令 \(z_1=-1\)，可得
    \(\bs{m}=\left( 2,\frac{1}{1-\lambda},-1 \right)\).

设平面 \(CFE\) 的法向量为 \(\bs{n}=(x_2,y_2,z_2)\)，则
\examdisplaycases{}{
\bs{n}\cdot\overrightarrow{FE}=-2y_2=0,\\
\bs{n}\cdot\overrightarrow{CF}=x_2+2z_2=0.
}
令 \(z_2=-1\)，可得
    \(\bs{n}=(2,0,-1)\).

从而 \(\bs{m}\cdot\bs{n}=5\)，\(\left| \bs{m} \right|=\sqrt{5+\left( \frac{1}{1-\lambda} \right)^2}\)，\(\left| \bs{n} \right|=\sqrt5\).
又二面角 \(M-CF-E\) 的余弦值为 \(\frac{\sqrt5}{3}\)，所以 \(\cos\langle \bs{m},\bs{n} \rangle=\frac{\bs{m}\cdot\bs{n}}{\left| \bs{m} \right|\left| \bs{n} \right|}=\frac{5}{\sqrt{5+\left( \frac{1}{1-\lambda} \right)^2}\times\sqrt5}=\frac{\sqrt5}{3}\).
整理可得 \((\lambda-1)^2=\frac{1}{4}\).
故 \(\lambda=\frac{1}{2}\)（\(\lambda=\frac{3}{2}\) 舍去）.

因此
    \(\frac{A_1M}{A_1B_1}=\frac{1}{2}\).
\end{solution}

\begin{problem}
为加快新冠肺炎检测效率，某检测机构采取“\(k\) 合 1 检测法”，即将 \(k\) 个人的拭子样本合并检测，若为阴性，则可以确定所有样本都是阴性的；若为阳性，则还需要对本组的每个人再做检测. 现有 \(100\) 人，已知其中 \(2\) 人感染病毒.

\begin{enumerate}
    \item
    \circled{1}若采用“10 合 1 检测法”，且两名患者同一组，求总检测次数；

    \circled{2}已知 \(10\) 人分成一组，分 \(10\) 组，两名感染患者在同一组的概率为 \(\frac{1}{11}\)，定义随机变量 \(X\) 为总检测次数，求检测次数 \(X\) 的分布列和数学期望 \(E(X)\)；
    \item 若采用“5 合 1 检测法”，检测次数 \(Y\) 的期望为 \(E(Y)\)，试比较 \(E(X)\) 和 \(E(Y)\) 的大小（直接写出结果）.
\end{enumerate}

\end{problem}

\begin{solution}
（1）\circled{1}对每组进行检测，需要 \(10\) 次；再对结果为阳性的组每个人进行检测，需要 \(10\) 次；

所以总检测次数为 \(20\) 次.

\circled{2}由题意，\(X\) 可以取 \(20\)，\(30\)，\(P(X=20)=\frac{1}{11}\)，\(P(X=30)=1-\frac{1}{11}=\frac{10}{11}\).
则 \(X\) 的分布列为：

\begin{center}
\begin{tabular}{|c|c|c|}
\hline
\(X\) & 20 & 30 \\
\hline
\(P\) & \(\frac{1}{11}\) & \(\frac{10}{11}\) \\
\hline
\end{tabular}
\end{center}

所以 \(E(X)=20\times\frac{1}{11}+30\times\frac{10}{11}=\frac{320}{11}\).

（2）由题意，\(Y\) 可以取 \(25\)，\(30\). 设两名感染者在同一组的概率为 \(p\)，则
    \(P(Y=25)=p\)，\(P(Y=30)=1-p\).
所以
所以 \(E(Y)=25p+30(1-p)=30-5p\).
于是：若 \(p=\frac{2}{11}\)，则 \(E(X)=E(Y)\)；

若 \(p>\frac{2}{11}\)，则 \(E(X)>E(Y)\)；

若 \(p<\frac{2}{11}\)，则 \(E(X)<E(Y)\).
\end{solution}

\begin{problem}
已知函数 \(f(x)=\frac{3-2x}{x^2+a}\).

\begin{enumerate}
    \item 若 \(a=0\)，求 \(y=f(x)\) 在 \((1,f(1))\) 处的切线方程；
    \item 若函数 \(f(x)\) 在 \(x=-1\) 处取得极值，求 \(f(x)\) 的单调区间，以及最大值和最小值.
\end{enumerate}

\end{problem}

\begin{solution}
（1）当 \(a=0\) 时，
    \(f(x)=\frac{3-2x}{x^2}\)，
则
    \(f'(x)=\frac{2(x-3)}{x^3}\).
所以
    \(f(1)=1\)，\(f'(1)=-4\).
此时，曲线 \(y=f(x)\) 在点 \((1,f(1))\) 处的切线方程为
    \(y-1=-4(x-1)\)，
即
    \(4x+y-5=0\).

（2）因为 \(f(x)=\frac{3-2x}{x^2+a}\)，则 \(f'(x)=\frac{2(x^2-3x-a)}{(x^2+a)^2}\).
由题意可得 \(f'(-1)=\frac{2(4-a)}{(a+1)^2}=0\)，
解得 \(a=4\).
于是 \(f(x)=\frac{3-2x}{x^2+4}\)，\(f'(x)=\frac{2(x+1)(x-4)}{(x^2+4)^2}\).
由此可知，函数 \(f(x)\) 的增区间为 \((-\infty,-1)\)，\((4,+\infty)\)，单调递减区间为 \((-1,4)\).
当 \(x<\frac{3}{2}\) 时，\(f(x)>0\)；当 \(x>\frac{3}{2}\) 时，\(f(x)<0\).
所以 \(f(x)_{\max}=f(-1)=1\)，\(f(x)_{\min}=f(4)=-\frac{1}{4}\).
\end{solution}

\begin{problem}
已知椭圆 \(E:\frac{x^2}{a^2}+\frac{y^2}{b^2}=1\)（\(a>b>0\)）过点 \(A(0,-2)\)，以四个顶点围成的四边形面积为 \(4\sqrt5\).

\begin{enumerate}
    \item 求椭圆 \(E\) 的标准方程；
    \item 过点 \(P(0,-3)\) 的直线 \(l\) 斜率为 \(k\)，交椭圆 \(E\) 于不同的两点 \(B\)，\(C\)，直线 \(AB\)，\(AC\) 交 \(y=-3\) 于点 \(M\)，\(N\)，若 \(\left| PM \right|+\left| PN \right|\le15\)，求 \(k\) 的取值范围.
\end{enumerate}

\end{problem}

\begin{solution}
（1）因为椭圆过 \(A(0,-2)\)，故 \(b=2\).
因为四个顶点围成的四边形的面积为 \(4\sqrt5\)，故 \(\frac{1}{2}\times2a\times2b=4\sqrt5\)，
解得 \(a=\sqrt5\).
故椭圆的标准方程为 \(\frac{x^2}{5}+\frac{y^2}{4}=1\).

（2）设 \(B(x_1,y_1)\)，\(C(x_2,y_2)\).

\bitmapfigure[max width=5.6cm,max totalheight=4.8cm]{img/2021/beijing/q20_solution_fig1.png}

因为直线 \(BC\) 的斜率存在，故 \(x_1x_2\ne0\).

故直线 \(AB\) 的方程为
    \(y=\frac{y_1+2}{x_1}x-2\).
令 \(y=-3\)，则
    \(x_M=-\frac{x_1}{y_1+2}\).

直线 \(AC\) 的方程为
    \(y=\frac{y_2+2}{x_2}x-2\).
令 \(y=-3\)，则
    \(x_N=-\frac{x_2}{y_2+2}\).

直线 \(BC\) 的方程为 \(y=kx-3\).
由
    \examdisplaycases{}{
    y=kx-3,\\
    4x^2+5y^2=20
    }
可得 \((4+5k^2)x^2-30kx+25=0\).
故 \(\Delta=900k^2-100(4+5k^2)>0\)，解得 \(k<-1\)或\(k>1\).

又 \(x_1+x_2=\frac{30k}{4+5k^2}\)，\(x_1x_2=\frac{25}{4+5k^2}\)，
故 \(x_1x_2>0\)，所以 \(x_Mx_N>0\).

从而 \(\left| PM \right|+\left| PN \right|=\left| x_M+x_N \right|=\left| \frac{x_1}{y_1+2}+\frac{x_2}{y_2+2} \right|\).
又因为 \(y_1=kx_1-3\)，\(y_2=kx_2-3\)，所以 \(\left| PM \right|+\left| PN \right|=\left| \frac{x_1}{kx_1-1}+\frac{x_2}{kx_2-1} \right|=\left| \frac{2kx_1x_2-(x_1+x_2)}{k^2x_1x_2-k(x_1+x_2)+1} \right|\).
代入 \(x_1+x_2\)，\(x_1x_2\) 可得 \(\left| PM \right|+\left| PN \right|=5|k|\).

由 \(\left| PM \right|+\left| PN \right|\le15\)，得 \(5|k|\le15\)，即 \(|k|\le3\).

综上，
    \(-3\le k<-1\)或\(1<k\le3\).
\end{solution}

\begin{problem}
定义 \(R_p\) 数列 \(\{a_n\}\)：对实数 \(p\)，满足：
\begin{circlelist}
\item \(a_1+p\ge0\)，\(a_2+p=0\)；
\item \(\forall n\in\N^*\)，\(a_{4n-1}<a_{4n}\)；
\item \(a_{m+n}\in\{a_m+a_n+p,\ a_m+a_n+p+1\}\)，\(m,n\in\N^*\).
\end{circlelist}

\begin{enumerate}
    \item 对于前 \(4\) 项 \(2\)，\(-2\)，\(0\)，\(1\) 的数列，可以是 \(R_2\) 数列吗？说明理由；
    \item 若 \(\{a_n\}\) 是 \(R_0\) 数列，求 \(a_5\) 值；
    \item 是否存在 \(p\)，使得存在 \(R_p\) 数列 \(\{a_n\}\)，对 \(\forall n\in\N^*\)，\(S_n\ge S_{10}\)？若存在，求出所有这样的 \(p\)；若不存在，说明理由.
\end{enumerate}

\end{problem}

\begin{solution}
（1）由性质 \circled{3} 结合题意可知
    \(0=a_3\in\{a_1+a_2+2,\ a_1+a_2+2+1\}=\{2,3\}\)，
矛盾，故前 \(4\) 项 \(2\)，\(-2\)，\(0\)，\(1\) 的数列，不可以是 \(R_2\) 数列.

（2）由性质 \circled{1} 可得 \(a_1\ge0\)，\(a_2=0\).
由性质 \circled{3} 可知 \(a_{m+2}\in\{a_m,\ a_m+1\}\)，
因此 \(a_3=a_1\ \text{或}\ a_3=a_1+1\)，\(a_4=0\ \text{或}\ 1\).

若 \(a_4=0\)，由性质 \circled{2} 可知 \(a_3<a_4\)，即 \(a_1<0\) 或 \(a_1+1<0\)，矛盾.

若 \(a_4=1\)，\(a_3=a_1+1\)，则由 \(a_3<a_4\) 有 \(a_1+1<1\)，矛盾.

因此只能是 \(a_4=1\)，\(a_3=a_1\).

又因为 \(a_4=a_1+a_3\ \text{或}\ a_4=a_1+a_3+1\)，所以 \(1=2a_1\ \text{或}\ 1=2a_1+1\)，
即 \(a_1=\frac{1}{2}\ \text{或}\ a_1=0\).

若 \(a_1=\frac{1}{2}\)，则 \(a_2\in\{a_1+a_1,\ a_1+a_1+1\}=\{1,2\}\)，
这与 \(a_2=0\) 矛盾，故只能有 \(a_1=0\).

于是 \(\{a_n\}\) 的前四项为：\(0,0,0,1\).

下面用归纳法证明：\(a_{4n+i}=n\ (i=1,2,3)\)，\(a_{4n+4}=n+1\)，\(n\in\N\).
当 \(n=0\) 时，经验证命题成立. 假设当 \(n\le k\) 时命题成立，则当 \(n=k+1\) 时，
    \(a_{4k+5}\in\{a_j+a_{4k+5-j},\ a_j+a_{4k+5-j}+1\mid 1\le j\le4k+4\}\)，
由归纳假设可知这些值只可能是 \(k\) 或 \(k+1\). 若 \(a_{4k+5}=k\)，取 \(k=0\) 可得 \(a_5=0\)，而由
    \(a_5\in\{a_1+a_4,\ a_1+a_4+1\}=\{1,2\}\)，
矛盾，所以 \(a_{4k+5}=k+1\).

由性质 \circled{3} 可得 \(a_{4k+6}\in\{a_1+a_{4k+5},\ a_1+a_{4k+5}+1\}=\{k+1,k+2\}\)，又 \(a_{4k+6}\in\{a_3+a_{4k+3},\ a_3+a_{4k+3}+1\}=\{k,k+1\}\)，故 \(a_{4k+6}=k+1\).

又由 \(a_{4k+8}\in\{a_2+a_{4k+6},\ a_2+a_{4k+6}+1\}=\{k+1,k+2\}\)，且 \(a_{4k+8}\in\{a_4+a_{4k+4},\ a_4+a_{4k+4}+1\}=\{k+2,k+3\}\)，故 \(a_{4k+8}=k+2\).

再由
    \(a_{4k+7}\in\{a_1+a_{4k+6},\ a_1+a_{4k+6}+1\}=\{k+1,k+2\}\)，
又因为 \(a_{4k+7}<a_{4k+8}\)，故
    \(a_{4k+7}=k+1\).
于是命题成立.

综上可得
    \(a_5=a_{4\times1+1}=1\).

（3）令
    \(b_n=a_n+p\).
由性质 \circled{3} 可知，对任意 \(m,n\in\N^*\)，
    \(b_{m+n}=a_{m+n}+p\in\{a_m+p+a_n+p,\ a_m+p+a_n+p+1\}=\{b_m+b_n,\ b_m+b_n+1\}\).
又因为
    \(b_1=a_1+p\ge0\)，\(b_2=a_2+p=0\)，\(b_{4n-1}=a_{4n-1}+p<a_{4n}+p=b_{4n}\)，
所以数列 \(\{b_n\}\) 为 \(R_0\) 数列.

由（2）可知
    \(b_{4n+i}=n\ (i=1,2,3)\)，\(b_{4n+4}=n+1\).
因此
    \(a_{4n+i}=n-p\ (i=1,2,3)\)，\(a_{4n+4}=n+1-p\).

由
    \(S_{11}-S_{10}=a_{11}=a_{4\times2+3}=2-p\ge0\)
和
    \(S_9-S_{10}=-a_{10}=-a_{4\times2+2}=-(2-p)\ge0\)
可得
    \(p=2\).

当 \(p=2\) 时，\(a_1,a_2,\cdots,a_{10}\le0\)，而 \(a_j\ge0\)（\(j\ge11\)），故对 \(\forall n\in\N^*\)，均有
    \(S_n\ge S_{10}\).

因此，存在这样的 \(p\)，且
    \(p=2\).
\end{solution}
